% ===================================================================
%  圖表第 9 號 — 山。森林。
%
%  Traduction, faite en 2026, en cantonais ecrit d'aujourd'hui, en
%  caracteres traditionnels, sur l'ido de texto/io. Regles de la
%  colonne : voir 10-toubiu-01.tex. Deux scenes, deux numerotations :
%  les cles portent c1 et c2.
%
%  LA TROISIEME TRANSCRIPTION DE LA COLONNE EST ICI, ET C'EST LA PLUS
%  LONGUE. Apres 芝士 (tableau 7) et 車厘子 (tableau 8) vient la
%  fraise :
%
%    士多啤梨 (11)   la fraise    Pekin : 草莓
%
%  Quatre caracteres pour « strawberry », un par syllabe cantonaise,
%  et pas un des quatre qui veuille dire quoi que ce soit ici. Pekin
%  a compose « baie des herbes » ; Hong Kong a ecoute le mot. La
%  regle degagee au tableau 8 tient sur les trois cas.
%
%  ET LES PETITES BETES DE LA FORET SE NOMMENT AUTREMENT QU'AU NORD :
%
%    烏蠅 (3)   la mouche     Pekin : 苍蝇
%    蟻竇 (6)   la fourmiliere Pekin : 蚁穴
%
%  « 蟻竇 » est le troisieme 竇 de la colonne, apres 雞竇 et 狗竇 au
%  tableau 7 : c'est le mot cantonais du gite, et il loge la fourmi
%  comme il logeait la poule.
%
%  LE PROMENEUR S'HABILLE EN CANTONAIS :
%
%    皮褸 (42)  la pelisse   Pekin : 皮袄   — 褸 est le manteau, et
%                                            le nord ne l'ecrit pas
%    背囊 (47)  le sac a dos Pekin : 背包
%    水浸 (44)  l'inondation Pekin : 洪水
%
%  LES METIERS EN 佬 PASSENT LA VINGTAINE : 礦佬, 打石佬,
%  伐木佬 (33), 燒炭佬 (41), 偷獵佬 (22). Et le cerf (50) prend le
%  suffixe 公 du tableau 7 — 鹿公, comme 雞公 et 羊公.
%
%  QUATRE BLOCS ONT DEMANDE L'INVERSION PREVENTIVE, tous sur le meme
%  patron « l'outil de l'ouvrier » : la scie du bucheron et la hache
%  du journalier (c2-07-1), le cor des chasseurs a cheval (c2-09-1),
%  le chien du garde (c2-05-1). Le chinois aurait nomme l'homme
%  d'abord ; on nomme l'outil, puis on dit a qui il est.
% ===================================================================

%%K t09-apar-1 apar
\VUsaut{4.0mm}
\VUtitre{15.0pt}{60}{圖表第 9 號}
\VUsaut{3.0mm}

%%K t09-tit-1 sub
\VUcentre{12.0pt}{0}{\textit{第一場。}}
\VUsaut{3.0mm}

%%K t09-tit-2 sub
\VUpk{11.4pt}{40}{山。--- 溫泉鎮。}
\VUsaut{3.0mm}

%%K t09-c1-01-1 p
1. --- 我嚟咗普拉茲八日嘞。當日我企喺庇里牛斯山嗰道靚到極嘅\VUgras{山脈}
\textsuperscript{(1)} 前面，心入面嗰份佩服，我真係講唔出：佢啲\VUgras{山脊}
\textsuperscript{(2)} 喺藍色天空上面，好似一條巨型花邊噉。

%%K t09-c1-02-1 p
2. --- 我從來估唔到庇里牛斯啲\VUgras{山峰} \textsuperscript{(3)} 有咁高
\VUgras{海拔；}我望住嗰啲靚到極嘅\VUgras{冰川} \textsuperscript{(5)}，同埋啲蓋住雪嘅
\VUgras{山尖} \textsuperscript{(4)}——嗰度成日滾落好嚇人嘅\VUgras{雪崩}——真係睇到呆咗。

%%K t09-tit-3 sub
\VUsaut{3.0mm}
\VUpk{10.2pt}{40}{山景。}
\VUsaut{3.0mm}

%%K t09-c1-03-1 p
3. --- 到埗第二日，我就去咗普拉茲附近嗰座熄咗火嘅老\VUgras{火山}
\textsuperscript{(6)}。個\VUgras{火山口}邊又乾又光禿禿，仲清清楚楚睇到當年
\VUgras{熔岩}流過嘅痕跡——幾百年前流落\VUgras{山坡} \textsuperscript{(7)} 嗰啲。我喺其中一道
\VUgras{斜坡}度，仲搵到幾嚿嗰啲熔岩。

%%K t09-c1-04-1 p
4. --- 普拉茲喺邊境上面。有座\VUgras{堡壘} \textsuperscript{(8)} 守住個
\VUgras{隘口} \textsuperscript{(27)}，將法國同西班牙分開；佢令人諗返起萊茵河邊嗰啲
舊城堡：佢哋企喺石頂上面，好似盯實\VUgras{獵物}噉。

%%K t09-c1-05-1 p
5. --- 有隻大\VUgras{鷹} \textsuperscript{(9)}，個巢就喺\VUgras{炮台} \textsuperscript{(8)} 唔遠處，
成日引到我望。

%%K t09-c1-05-1 p suite
呢隻\VUgras{猛禽}把\VUgras{嘴}好夠力，佢時時用對\VUgras{爪}捉住隻羊仔或者山羊仔
返去餵佢啲細嘅。佢對\VUgras{翼}大到，就算孭住咁重嘅嘢，都可以滑翔好耐。

%%K t09-tit-4 sub
\VUsaut{3.0mm}
\VUpk{10.2pt}{40}{行山客。}
\VUsaut{3.0mm}

%%K t09-c1-06-1 p
6. --- 呢頭最舒服嘅散步，就係行去「高」嗰道\VUgras{瀑布} \textsuperscript{(10)}。度
\VUgras{飛瀑}打住漩渦跌落嚟，跟住化成一條靚\VUgras{河仔，}冇入城西邊嗰片
\VUgras{杉樹林} \textsuperscript{(11)}。我哋穿過呢片樹林一路上，去到個\VUgras{氣象台}
\textsuperscript{(12)}；企喺個\VUgras{觀景台}頂，成區嘅\VUgras{全景}都收晒喺眼底。
啲\VUgras{風景}靚到唔講得，\VUgras{大自然}好似將自己有嘅寶都倒晒畀呢笪地方。

%%K t09-c1-07-1 p
7. --- 落山我哋搭\VUgras{纜車} \textsuperscript{(13)}，喺個細站停低，就喺一座舊
\VUgras{封建城堡}嘅\VUgras{頹垣} \textsuperscript{(14)} 隔籬——以前普拉茲啲領主住嗰座。

%%K t09-c1-08-1 p
8. --- 可憐嘅城堡！佢望住山下面嗰個靚靚哋嘅\VUgras{溫泉鎮} \textsuperscript{(15)}
——而家已經係間時髦嘅\VUgras{溫泉站}——唔知佢心入面點諗呢？

%%K t09-c1-08-2 p
呢度\VUgras{氣候}咁舒服，度\VUgras{礦泉水}又咁見效，就算喺間\VUgras{療養院}
\textsuperscript{(17)} 度\VUgras{療養，}都真係一件樂事。

%%K t09-tit-5 sub
\VUsaut{3.0mm}
\VUpk{10.2pt}{40}{舒舒服服嘅溫泉鎮。}
\VUsaut{3.0mm}

%%K t09-c1-09-1 p
9. --- 何況呢度乜都做齊晒，令人住得舒服。起咗道大\VUgras{高架橋}
\textsuperscript{(16)}，等火車入得嚟；溫泉院嗰個\VUgras{公園} \textsuperscript{(18)}——入面有
\VUgras{音樂會}聽——就佈置得好有味道。起咗好幾間靚靚哋嘅
「\VUgras{別墅}」\textsuperscript{(19)}，圍住個\VUgras{賭場} \textsuperscript{(21)}；
仲有條保養得好好嘅\VUgras{大馬路}
\textsuperscript{(22)}，來往就方便好多。

%%K t09-c1-10-1 p
10. --- 一道普普通通嘅\VUgras{斜坡} \textsuperscript{(23)}，就將條\VUgras{路}
\textsuperscript{(20)} 同真真正正嘅\VUgras{山谷} \textsuperscript{(25)} 隔開；谷底躺住個
靚靚哋嘅細\VUgras{湖} \textsuperscript{(26)}，供住一條清見底嘅\VUgras{河仔} \textsuperscript{(24)}。

%%K t09-c1-11-1 p
11. --- 山谷另一邊，有個\VUgras{鐵礦} \textsuperscript{(28)} 開緊。我去過幾次睇啲
\VUgras{礦佬}落\VUgras{豎井，}仲入過佢哋成日喺度嘅\VUgras{巷道，}睇佢哋鑿出啲
含住\VUgras{金屬}嘅\VUgras{礦石。}

%%K t09-c1-12-1 p
12. --- 礦場隔籬起咗間好大嘅\VUgras{鑄鐵廠，}即刻用得着礦入面挖出嚟嘅嘢。個
\VUgras{高爐} \textsuperscript{(29)} 燒嘅係呢頭多到不得了嘅\VUgras{煤，}又快又平就出到
\VUgras{生鐵。}

%%K t09-c1-13-1 p
13. --- 離礦場唔遠，仲開緊個\VUgras{石礦場} \textsuperscript{(30)}。老\VUgras{打石佬}
用把\VUgras{鶴嘴鋤}鑿出嚟嘅，唔係\VUgras{花崗石，}唔係\VUgras{斑岩，}連\VUgras{沙石}
都唔係，係起屋用嘅普通\VUgras{方石。}

%%K t09-c1-14-1 p
14. --- 呢笪地方最靚嘅裝飾之一，就係啲\VUgras{杉樹} \textsuperscript{(31)}。無論
\VUgras{峽谷} \textsuperscript{(32)} 谷底、\VUgras{急流} \textsuperscript{(33)} 岸邊、\VUgras{深淵}
\textsuperscript{(34)} 或者崖邊，佢哋啲幼針葉都染出一片綠。

%%K t09-tit-6 sub
\VUsaut{3.0mm}
\VUpk{10.2pt}{40}{行路。}
\VUsaut{3.0mm}

%%K t09-c1-15-1 p
15. --- 我趁住假期\VUgras{行遠啲。}啲喺山上面彎彎曲曲嘅\VUgras{路}
\textsuperscript{(35)}，行起上嚟唔使理幾遠，一行就幾個\VUgras{公里，}成日仲行到
成\VUgras{幾十公里。}

%%K t09-c1-15-2 p
啲路有\VUgras{界石} \textsuperscript{(36)} 同\VUgras{路標柱} \textsuperscript{(37)} 做記號；路上時時
撞到後生嘅\VUgras{山裏人} \textsuperscript{(38)}，腰間掛住個\VUgras{褡褳} \textsuperscript{(40)}，
孭住隻\VUgras{土撥鼠} \textsuperscript{(39)}；又或者撞到幾個\VUgras{吉卜賽人}
\textsuperscript{(41)}，佢件\VUgras{皮褸} \textsuperscript{(42)} 同佢頂爛晒嘅舊\VUgras{兜帽}
\textsuperscript{(43)} 襯到好古怪。佢用條鏈拖住隻馴咗嘅\VUgras{熊} \textsuperscript{(44)}。

%%K t09-tit-7 sub
\VUpk{10.2pt}{40}{爬山。}
\VUsaut{3.0mm}

%%K t09-c1-16-1 p
16. --- 我差唔多日日都跟住個\VUgras{嚮導} \textsuperscript{(48)} 去練\VUgras{爬山；}
孭住個\VUgras{背囊} \textsuperscript{(47)}、手攞住支「\VUgras{冰鎬}」，好似個真正
\VUgras{遊客} \textsuperscript{(46)} 噉衝上啲\VUgras{石堆} \textsuperscript{(45)} 嗰陣，我真係威到爆。

%%K t09-c1-17-1 p
17. --- 企喺山頂望落去，平原上面啲人望落唔大過隻螞蟻；一群\VUgras{羊}
\textsuperscript{(50)} 喺個\VUgras{牧場} \textsuperscript{(49)} 度食草，望落好似細路仔嘅玩具。
一棵\VUgras{松樹} \textsuperscript{(51)}，或者一間\VUgras{木屋} \textsuperscript{(52)}，喺片綠色上面
就只不過係一點斑。

%%K t09-c1-18-1 p
18. --- 行呢啲山嗰陣，我哋成日喺條\VUgras{山路仔} \textsuperscript{(53)} 度撞到
\VUgras{獵岩羚嘅} \textsuperscript{(54)}，佢哋膊頭孭住啱啱打低嗰隻嘢。

%%K t09-c1-19-1 p
19. --- 我本標本簿度夾咗一枝好值得留嘅\VUgras{蕨}
\textsuperscript{(55)}，同一枝靚粉紅色嘅\VUgras{石南花} \textsuperscript{(57)}——
兩枝都係我最後嗰次散步喺個細\VUgras{石洞} \textsuperscript{(56)} 門口摘嘅。

%%K t09-tit-8 sub
\VUsaut{3.0mm}
\VUcentre{12.0pt}{0}{\textit{第二場。}}
\VUsaut{3.0mm}

%%K t09-tit-9 sub
\VUpk{11.4pt}{40}{森林。--- 打獵。}
\VUsaut{3.0mm}

%%K t09-c2-01-1 p
1. --- 尋日我喺森林度過咗好正嘅一日。住喺入面嗰啲動物同\VUgras{昆蟲}
\textsuperscript{(5)} 幾咁勤力，睇到我好有興趣。

%%K t09-c2-01-2 p
隻\VUgras{蜘蛛} \textsuperscript{(1)} 喺兩枝樹枝之間織住個\VUgras{網} \textsuperscript{(2)}，
等隻飛過嘅\VUgras{烏蠅} \textsuperscript{(3)} 撞落嚟；隻\VUgras{蝸牛} \textsuperscript{(4)}
辛辛苦苦孭住個殼；隻鍬形蟲喺度巡佢嘅收穫；勤力嘅螞蟻喺\VUgras{蟻竇}
\textsuperscript{(6)} 隔籬忙到飛起——呢班細嘢行來行去、做嚟做去，勤力到唔尋常。

%%K t09-c2-02-1 p
2. --- 有條\VUgras{蛇} \textsuperscript{(7)}，我一眼望落以為係\VUgras{蝰蛇，}原來
淨係條普通\VUgras{草蛇；}佢匿咗喺樹頭下面，時不時伸出個幼幼嘅頭，成副樣好似
隨時想咬人噉。我面前有隻大\VUgras{鼻涕蟲} \textsuperscript{(8)} 慢慢噉爬，佢啱啱食咗
一粒呢度周圍都係嘅香\VUgras{士多啤梨} \textsuperscript{(11)}。

%%K t09-c2-03-1 p
3. --- 地面鋪滿\VUgras{林裏嘅花：}\VUgras{報春花} \textsuperscript{(9)}、
\VUgras{長春花} \textsuperscript{(10)}，仲有好多我唔識嘅草，喺一叢叢\VUgras{草}
\textsuperscript{(12)} 中間開住。

%%K t09-c2-04-1 p
4. --- 喺呢一頭，\VUgras{松露}差唔多同\VUgras{蘑菇} \textsuperscript{(14)} 一樣咁常見；
有隻護林員養嘅\VUgras{豬仔} \textsuperscript{(13)} 就饞嘴噉周圍拱，睇下搵唔搵到幾隻。

%%K t09-tit-10 sub
\VUsaut{3.0mm}
\VUpk{10.2pt}{40}{偷獵。}
\VUsaut{3.0mm}

%%K t09-c2-05-1 p
5. --- 我睇到一幕戲嘅\VUgras{結尾，}真係幾好笑。靠住隻\VUgras{矮腳獵犬}
\textsuperscript{(15)} 嘅鼻，\VUgras{護林員} \textsuperscript{(16)} 啱啱撞破咗個\VUgras{偷獵佬}
\textsuperscript{(22)}——就喺嗰條友準備孭走佢打低嘅\VUgras{野味} \textsuperscript{(17)} 嗰陣。

%%K t09-c2-05-1 p suite
偷獵佬情願唔要啲收穫，都唔想畀人拉，於是就走咗人；地上啲草度就剩返隻
\VUgras{野兔} \textsuperscript{(18)}、隻\VUgras{鹿乸} \textsuperscript{(19)}、隻\VUgras{麅}
\textsuperscript{(20)} 同隻\VUgras{野豬} \textsuperscript{(21)}，睇到個護林員笑騎騎。

%%K t09-c2-06-1 p
6. --- 森林入面啲樹之多，真係令人咋舌。\VUgras{山毛櫸} \textsuperscript{(23)} 同
\VUgras{白樺} \textsuperscript{(26)}、出\VUgras{松脂} \textsuperscript{(27)} 嘅\VUgras{松樹、}
\VUgras{橡樹} \textsuperscript{(29)}，仲有好多好多，啲樹椏都纏埋一齊。

%%K t09-c2-06-2 p
\VUgras{常春藤} \textsuperscript{(24)} 攀住啲高樹嘅樹幹，將成片\VUgras{樹林}
\textsuperscript{(25)} 染成閃住光嘅綠色；呢種綠同啲\VUgras{青苔} \textsuperscript{(31)}
淺啲、白啲嘅綠襯得好夾，隻\VUgras{啄木鳥} \textsuperscript{(30)} 就喺青苔度搵蟲食。

%%K t09-tit-11 sub
\VUsaut{3.0mm}
\VUpk{10.2pt}{40}{森林嘅景。}
\VUsaut{3.0mm}

%%K t09-c2-07-1 p
7. --- 啲老橡樹迷到我。佢哋啲粗\VUgras{樹根} \textsuperscript{(32)} 扭曲晒，一半
拱咗出泥面，望落夠曬硬淨夠曬有力；我真係唔明個\VUgras{業主} \textsuperscript{(35)}
點捨得斬咗佢哋，交畀把\VUgras{鋸} \textsuperscript{(34)}——即係\VUgras{伐木佬}
\textsuperscript{(33)} 嗰把。啲可憐嘅\VUgras{樹幹} \textsuperscript{(36)} 畀人剝清\VUgras{樹皮}
\textsuperscript{(37)}，又畀把\VUgras{斧頭} \textsuperscript{(38)} 劈開——\VUgras{散工}
\textsuperscript{(39)} 手上攞嗰把——睇到我心都實。

%%K t09-c2-08-1 p
8. --- 隔籬有個老\VUgras{燒炭佬} \textsuperscript{(41)}，用把鏟堆起一\VUgras{堆}
\textsuperscript{(42)} 炭；仲有個阿婆孭住一\VUgras{紮} \textsuperscript{(40)} \VUgras{枯柴，}係
斬低嘅樹上面啲細枝。

%%K t09-tit-12 sub
\VUsaut{3.0mm}
\VUpk{10.2pt}{40}{追獵。}
\VUsaut{3.0mm}

%%K t09-c2-09-1 p
9. --- 我本來想入\VUgras{護林員屋} \textsuperscript{(46)} 度歇一陣，但係行到個
\VUgras{湖仔} \textsuperscript{(43)} 附近——湖面上游住隻靚\VUgras{天鵝} \textsuperscript{(45)}
——先發覺頭先場\VUgras{水浸} \textsuperscript{(44)} 已經將條路變成一片\VUgras{沼澤。}
我唯有兜路走；行經森林邊嗰棵\VUgras{香柏} \textsuperscript{(49)}、幾棵\VUgras{白楊}
\textsuperscript{(48)} 同棵\VUgras{榆樹} \textsuperscript{(47)} 嗰陣，我見到由片\VUgras{矮樹林}
\textsuperscript{(54)} 度衝出一隻靚到極嘅\VUgras{鹿公} \textsuperscript{(50)}，後面
一群死咬唔放嘅\VUgras{獵狗} \textsuperscript{(51)} 追住；仲有把\VUgras{號角} \textsuperscript{(53)}
催住佢哋——係啲\VUgras{騎馬嘅獵人} \textsuperscript{(52)} 吹嘅。可憐嗰隻嘢已經冇晒力，
行到離我幾步遠就瞓低死咗。

%%K t09-c2-10-1 p
10. --- 護林員個仔畀\VUgras{追獵}嘅聲吸引咗出嚟，我哋兩個就一齊行返入森林。
佢頭先見到隻\VUgras{喜鵲} \textsuperscript{(56)} 企喺棵老橡樹嘅樹椏上面。佢估個巢
應該匿咗喺啲\VUgras{樹葉} \textsuperscript{(58)} 度，就想去掏。

%%K t09-c2-10-2 p
佢好似隻\VUgras{松鼠} \textsuperscript{(61)} 噉靈活，由呢枝\VUgras{樹椏} \textsuperscript{(55)}
爬去嗰枝；不過乜都搵唔到，就淨係嚇到隻老\VUgras{烏鴉}
\textsuperscript{(60)} 飛咗走——嗰隻本來企喺條大\VUgras{樹椏} \textsuperscript{(57)}
上面。

\VUsaut{4.0mm}
\VUfilet{22mm}
